%
%============= CHAPTER BREAK =================%
%
\chapter{[Jehudite-]Christian Neologisms}%
%
\vspace*{3mm}
%
\justifying

\noindent{}While this book is not focused on neologisms %
coined or first popularized by the author(s)/editor(s) of the texts in the New Testament, %
our research dug up such candidates. %
Many of these reveal close connections with other early Christian texts, %
with crossover at the lexical level between canonical and mid-2nd to mid-3rd century pseudepigraphical/apocryphal texts, %
and between canonical scriptures and known mid-2nd to mid-3rd century Christian writers. %
These crossovers challenge the common, baseless assumption %
that canonical texts must be earlier and non-canonical texts later, %
that authors/editors of non-canonical literature %
could not have played a part in the authoring/editing of canonical literature, %
and that Christian writers from the mid-2nd to mid-3rd century had no significant editorial role in the canonical NT texts. %

Various early Christian crossover neologisms are saved for later sections in this chapter. 
For these, we follow the organizing principle of sorting attestations according to the first witness outside the %
NT texts, whether that be an apocryphal, pseudepigraphical, or early Christian writer. %
When a term occurs in multiple writings that may be roughly contemporaneous, %
they are grouped under what is postulated %
as the earliest text/author outside of canonical candidates. %
Given the often uncertain dates and textual fluidity of their corresponding texts, %
these subsequent sections run in alphanumeric order by author or title, %
not by postulated chronology. %
Sometimes these crossover neologisms are provided in simple lists, %
and at other times briefly explored and analyzed.
If no apparent usage of a lemma or clear reception of its NT usage is found between the mid-2nd and mid-3rd century, %
then those lemma entries are placed under the general NT Sprung Neologisms list. \clearpage %

\section{NT Sprung Neologisms}

\noindent{}We begin with a list of the strongest candidates for genuine NT neologisms. %
These lack secure parallels before the mid-3rd century CE. %
The list is filtered according to the limited research scope described in the introduction, %
excluding named entities and Hebrew and Aramaic transliterations into Greek. %
The list provides any NT citation(s), %
then the next securely dated hit, based on a review of date-sorted TLG search results %
and filtering out late and/or dubiously dated/attributed hits. %
Neologisms found in Luke and the Pauline epistles %
always correspond to the canonical versions, but rarely to their Marcionite counterparts, %
with a few exceptions in Deutero-Pauline Ephesians (Marcionite Laodiceans) and Colossians. %
Exceptions in Marcionite Colossians are noted with a leading asterisk, thus *Col. \\[5mm] %
%
\begin{small} %
\noindent{}\textbf{αἰτίωμα} / \textit{aitiōma} / “legal charge” (Acts 25:7; \textit{Corp.Herm.} Stobaean fr. 23)\endnote{%
Nock 4:12, Litwa 2018:118 translates it here as “blame”.} \\[0.4mm] %
\noindent{}\textbf{ἀκατάκριτος} / \textit{akatakritos} / “uncondemned” (Acts 16:37, 22:25; \\ % 
 \makebox[3em][l]{} \textit{Acts.Chr.Pet.} 8.1)%
\endnote{\textit{The Acts of Christ and Peter} (ECCA 956, CANT 198), %
also known as \textit{Acta fabulosa} (BHG 1485f), %
which Snyder and Čéplö (2023:440--441) tentatively date between the 4th and 6th century CE. %
The term is clearly in evidence in late 4th century Greek works, most securely Chrysostom, \textit{ActH}.} \\[0.4mm] %
\noindent{}\textbf{ἀλλοτριεπίσκοψ}|\textbf{ἀλλοτριεπίσκοπος} / \textit{allotriepiskops}|\textit{allotriepiskopos} / \\ % 
 \makebox[3em][l]{} “anotherʼs overlooker” (1Pt 4:12; Epiph, \textit{Ancor}. 12.5) \\[0.4mm] % 
\noindent{}\textbf{ἀνεμίζομαι} / \textit{anemidzomai} / “to be wind-blown” (Jas 1:6; BasilC, \textit{ep.} 42.4) \\[0.4mm] %
\noindent{}\textbf{ἀνεύθετος} / \textit{aneuthetos} / “inconvenient” (Acts 27:12; JChrys, \textit{ActH} \\\makebox[3em][l]{}PG60:368--370) \\[0.4mm] %
\noindent{}\textbf{ἀπελεγμός} / \textit{apelegmos} / “refutation” (Acts 19:27; JChrys, \textit{ActH} PG60:297) \\[0.4mm] %
\noindent{}\textbf{ἀφθορία} / \textit{aphthoria} / “incorruption” (Tit 2:7; Themistius, \textit{Arist.phys.} 5.2:82) \\[0.4mm] %
\noindent{}\textbf{ἀφιλάγαθος} / \textit{aphilagathos} / “not-good-loving” (2Tim 3:3; BasilC, \textit{Reg.mor.} \\\makebox[3em][l]{}PG31:812) \\[1mm] %
\noindent{}\textbf{βολίζω} / \textit{bolidzō} / “to take soundings”, “to heave the lead”, %
“to probe|gauge \\\makebox[3em][l]{}the depth” (Acts 27:28; \textit{Acts.Chr.Pet.} 8.1)\endnote{%
As with the lemma ἀκατάκριτος, βολίζω also may well next appear securely in Chrysostomʼs late 4th century \textit{ActH} (PG60:369). %
In Paulʼs sea-voyage in NT Acts, the verb appears twice in quick succession: %
\vspace*{1mm} 
\begin{quote} Ὡς δὲ τεσσαρεσκαιδεκάτη νὺξ ἐγένετο διαφερομένων ἡμῶν ἐν τῷ Ἀδρίᾳ, %
κατὰ μέσον τῆς νυκτὸς ὑπενόουν οἱ ναῦται προσάγειν τινὰ αὐτοῖς χώραν. %
²⁸ καὶ \textbf{βολίσαντες} εὗρον ὀργυιὰς εἴκοσι, %
βραχὺ δὲ διαστήσαντες καὶ πάλιν \textbf{βολίσαντες} εὗρον ὀργυιὰς δεκαπέντε· \\[1mm]
Now on the fourteenth night it happened as we were passing through the Adriatic, %
at midnight the sailors suspected approaching some land for themselves. %
²⁸ And \textbf{depth-probing} they found twenty fathoms, %
but after observing a brief interval and again \textbf{depth-probing}, they found fifteen fathoms. %
\end{quote} %
\vspace*{1mm}
\noindent{}The description of measuring sea-floor depth seems ordinary enough, %
but the verb is evidently new, missing even from the extensive lexicographical works of Dysculos and Herodian, %
which sometimes covered domain-specific (including nautical) terminology. %

βολίζω apparently derives from βολίς / \textit{bolis} / “missile”, “javelin”, “sounding-lead” (LSJ), %
but may also closely relate to βόλος / \textit{bolos} / “cast”, “throw with a casting-net”, “catch” (LSJ). %
Traditional and more commonplace Greek verbs for assessing water depth at sea are καταπειράζω / \textit{katapeiradzō} / “to probe”, “to test depth” %
(e.g., Polybius, \textit{Hist}. 4.11.6), ἀναμετρέω / \textit{anametreō} / “to measure up” %
(Agatharchides, \textit{mar.Eryth.} 83; Posidonus, fr. 9), both often with ὄργυια / \textit{orguia} / “fathom”, “six feet”. 
Sometimes εὑρίσκω / \textit{heuriskō} / “to find” is used on its own for this purpose (Diodorus Siculus, \textit{Bibl.hist.} 3.40.3). %
The same unusual verb also appears in the \textit{Acts.Chr.Pet.} 8:1: %

\vspace*{1mm} 
 \begin{quote}Τριήμερον δὲ \textbf{βολίσαντες} ἀγκύρας ἑορτὴν ἀπετέλουν. %
 Ὁ δὲ Πέτρος ἠγαλλιᾶτο δοξάζων τὸν Θεόν. \\[1mm] %
 \textls[-10]{Now after three days of \textbf{depth-probing} by anchor, they were performing festival. %
 But Peter was rejoicing greatly, glorifying the God.} %
 \end{quote} % 
\vspace*{1mm} 

\noindent{}While the verb is strange, the usage seems ordinary enough, %
until the broader narrative unfolds. %
To Peter the “shipmaster” / ναύκληρος (Christ in disguise) offers his son (also Christ in disguise), %
a mere “lad” / μειράκιον, as a slave (8:2). %
Peter objects, since Jesus forbade him from %
carrying gold or silver, and thus he lacks the means to buy the lad (8:3). %
The shipmaster then offers Peter ten gold coins so that Peter can complete the transaction (8:4), %
persuading him to accept the monetary and human gift, %
and even providing a bill of sale (8:5--6). %
The bill stipulated that the lad, here named “Messias” / Μεσσίαν, %
would serve Peter until Peter died, and thereafter would be returned to his father. %
In view of the broader narrative, %
one wonders whether the use of the unusual word βολίζω / “depth-probing”, %
particularly in connection with a “festival”, %
might have disturbing connotations associated with enslavement. %
} \\[0.4mm] %
\noindent{}\textbf{γογγυστής} / \textit{gongustēs} / “grumbler” (Jude 1:16; Euseb, \textit{Isa.com.} 1.97) \\[0.4mm] %
\noindent{}\textbf{δεξιολάβος} / \textit{deksiolabos} / “spearman” (Acts 23:23; JChrys, \textit{ActH} \\\makebox[3em][l]{}PG60:340) \\[0.4mm] %
\noindent{}\textbf{δευτερόπρωτος} / \textit{deuteroprōtos} / “second-first” (Lk 6:1; JChrys, \textit{MtHom90} \\\makebox[3em][l]{}PG57:433) \\[0.4mm] %
\noindent{}\textbf{διακατελέγχομαι} / \textit{diakatelegchomai} / “to confute totally” (Acts 18:28; JChrys, \\\makebox[3em][l]{}\textit{ActH} PG60:283)\endnote{%
The lemma appears in one Greek papyrus (michigan.apis.1470, but it is a witness to Acts 18--21).} \\[0.4mm]
\noindent{}\textbf{δισμυριάς} / \textit{dismurias} / “the two-thousand” %
(Rev 9.16; Andreas, \textit{Apoc.} 9.27.9)\endnote{%
The noun form (δισμυριάς, ‑άδος, ἡ) is only used in Revelation and commentators who quote it. %
The adjective (δισμύριοι, ‑αι, ‑α; δισμύριος, α, ον) is sometimes substantival, %
and consistently used from the 5th c. BCE to the late Byzantine era.} \\[0.4mm] %
\noindent{}\textbf{ἐκπερισσῶς} / \textit{ekperissōs} / “more exceedingly” (Mk 14:31; Ps-Gregentius, \\\makebox[3em][l]{}\textit{Vit.Greg.} 7) \\ %
\noindent{}\textbf{ἐμπαιγμονή} / \textit{empaigmonē} / “mockery” (2Pt 3:3; CyrAlex, \textit{incarn.unig.})\endnote{Durand SC97:240, Aubert, p. 695.} \\ %
\noindent{}\textbf{ἐπιδιατάσσομαι} / \textit{epidiatassomai} / “to add to a will” (Gal 3:15; \\\makebox[3em][l]{}\textit{Adv.Cataphryg.} fr. 1)\endnote{\textit{Adv. Cataphrygas} is CPG1328 and TLG1131.001, with the digital text from Routh 1846:183.} \\[0.4mm] %
 \noindent{}\textbf{ἐπισυντρέχω} / \textit{episuntrechō} / “to run together” (Mk 9:25; Ps-Hippolytus \\\makebox[3em][l]{}of Rome, \textit{consum.mund.} 42) \\ %
\noindent{}\textbf{ἑτεροδιδασκαλέω} / \textit{heterodidaskaleō} / “to teach differently” (1Tim 1:3, 6:3; \\\makebox[3em][l]{}Ign7, \textit{Poly}. 3.1) \\[0.4mm] %
\noindent{}\textbf{εὐπροσωπέω} / \textit{euprosōpēsai} / “to present well”|“to look good” (Gal 6:12; \\\makebox[3em][l]{}Euseb, \textit{Gal.fr.} p. 52, JChrys, \textit{Gal.com.} PG61:677--678)\endnote{%
This term is unattested but restored for Marcionʼs \textit{Apostolos} in our TANZ edition, %
but should be removed in light of this finding.} \\[0.4mm] %
\noindent{}\textbf{καταλιθάζω} / \textit{katalithadzō} / “to stone to death” (Lk 20:6; Nys, \\\makebox[3em][l]{}\textit{Eunom}. 1.1.265) \\[0.4mm] %
\noindent{}\textbf{καταστρηνιάω} / \textit{katastrēniaō} / “to behave wantonly toward” (1Tim 5:11; \\\makebox[3em][l]{}Ign13, \textit{Ant}. 11.1) \\[0.4mm] %
\noindent{}\textbf{κατεφίσταμαι} / \textit{katephistamai} / “to rise up against” (Acts 18:12; JChrys, \\\makebox[3em][l]{}\textit{ActH} PG60:277) \\[0.4mm] %
\noindent{}\textbf{μετάλημψις} / \textit{metalēmpsis} / “a sharing” (1Tim 4:3; Theon, \textit{Ptol}. p. 879)\endnote{%
The Attic form, μετάληψις, lacked the medial μ, and remained commonplace throughout the Christian era. %
The form μετάλημψις appears in 1Tim in mss 01 02 etc, but is also missing from many early mss.} \\[0.4mm] %%
 \noindent{}\textbf{ὀλιγόπιστος} / \textit{oligopistos} / “mini-faithful” (*Ev 12:28, Mt 6:30, 8:26, 14:31, 16:8; %
\\\makebox[3em][l]{}ClemAlex, \textit{Paed}. 2.10b.103.1, \textit{Strom}. 7.42.4) \\[0.4mm] %
\noindent{}\textbf{ὀρθοποδέω} / \textit{orthopodeō} / “to walk straight” (*Gal 2:14; Orig, \textit{Mt.com.} 16.24) \\[0.4mm] %
\noindent{}\textls[-15]{\textbf{ὀχλοποιέω} / \textit{ochlopoieō} / “to make a riot” (Acts 17:5; JChrys, \textit{ActH} PG60:263)}\\[0.4mm] %
 \noindent{}\textbf{παραφρονία} / \textit{paraphronia} / “madness” (2Pt 2:16; JChrys, \textit{fr.ep.cath.} PG64:1060) \\[0.4mm] %
\noindent{}\textbf{πορφυρόπωλις}|\textbf{πορφῠροπώλης} / \textit{porphuropōlis}|\textit{porphuropōlēs} / “purple-dyer” \\\makebox[3em][l]{}(Acts 16:14; Epiph, \textit{Pan}. GCS 3:485)\endnote{%
LSJ cites IG Rom.4.1071 (Cos) for the feminine, yet notes that the masculine first declension πορφῠροπώλης / \textit{porphuropōlēs}, %
is attested in the inscription. See Humann 1898:121, no. 156. %
The inscription reads ἡ σορὸς Μ(άρκου) Αὐρ(ηλίου) Ἀλεξάνδρου Μοσχιανοῦ βουλευτοῦ πορφυροπώλου / %
“The sacrophagus of M(arkus) Aur(elius) Alexander Moschianos, council-member, purple-dyer”. 
This inscription is also classified as IHierapJ 156 — PH271772 — AGRW 10527 at \url{www.philipharland.com/greco-roman-associations/?p=10527}. %
See also IHerapP 37 — PH268802 — AGRW 10581 at \url{www.philipharland.com/greco-roman-associations/?p=10581}. %
Harland dates both inscriptions between 212--300 CE. %
The cognate lemma πορφυροπωλικὴ / \textit{porphuropōlikē} / “purple-dye craft” is %
attested by the grammarian (and tutor to emperor Lucius Verus) Harpocration, %
\textit{Lexic}., s.v., Ἁλουργοπωλική, p. 13. %
Recent scholarship, drawing mainly on biographical evidence found in Greek papyri, %
has Harpocration flourishing from the late 2nd to early 3rd century (Keaney 1991). %
A search of the common stem on papyri.info turns up 13 hits, but all of them are 4th century CE or later. %
While Harpocration is not the clear basis for the term in Acts, the confluence of evidence %
suggests a late 2nd or early 3rd century setting for this characterization of the Lydia character of Acts 16.} \\[0.4mm] %
\noindent{}\textbf{προενάρχομαι} / \textit{proenarchomai} / “to begin before” (2Cor 8:6, 10; JChrys, %
\\\makebox[3em][l]{}\textit{2Cor.hom.} PG61:515, 518, 529--530) \\[0.4mm]
\noindent{}\textbf{προκυρόω}|\textbf{προκυρόομαι} / \textit{prokuroō}|\textit{prokuroomai} / “to ratify before” %
\\\makebox[3em][l]{}(Gal 3:17; Ps-Orig, \textit{Ps.fr.} 77.29--30) \\[0.4mm] %
\noindent{}\textbf{προμαρτύρομαι} / \textit{promarturomai} / “to witness beforehand” (1Pet 1:11; %
\\\makebox[3em][l]{}Euseb, \textit{H.E.} 1.3.6) \\[0.4mm] %
\noindent{}\textbf{προσεάω} / \textit{proseaō} / “to suffer going further” (Acts 27:7; JChrys, \textit{Act.cat.})\endnote{%
See Cramer 1938, 3:397 for the text. %
LSJ notes that PLond.5.1790.7, 5th--6th CE, %
attests the lemma. It did not appear in any of my searches of digitized Greek papyri or inscriptions.} \\[0.4mm] %
\noindent{}\textbf{πρόσλημψις} / \textit{proslēmpsis} / “acquisition” (Rom 11:15; DmC, \textit{Eccl.com.} 7--8.8)\endnote{%
Kramer \& Krebber 1972:222. Previous instances of the lemma in the TLG, digitized Greek papyri, and Greek inscriptions omit the medial μ.} \\[0.4mm] %
\noindent{}\textbf{πυρράζω} / \textit{purradzō} / “to redden” (Mt 16:2--3, BasilC, \textit{Hexaem}. 6.4)\endnote{%
A search of the major database of Greek papyri turns up seven hits for this lemma, %
but these are all incorrectly tagged as verbs, when they all reflect the adjective πυρρός, ‑ά, ‑όν, or else the noun πύρρα, ‑ας, ἡ.} \\[0.4mm] %
\noindent{}\textls[-15]{\textbf{σιτομέτριον} / \textit{sitometrion} / “grain-measure” (*Ev 12:42; Orig, \textit{Jo.com.} 20.2.8)}\endnote{%
The lemma appears after 152 CE in an inscription in Lycia, Asia Minor: FdXanthos VII 67, Lycia, W. — Xanthos (Kınık) — Letoion (Bozoluk) — SEG 30.1535 — PH283274. %
This lemma, or a close cognate, also appears on an 2nd century CE inscription in Aegean Kos, IKosS EF  SEG 57 (2007), no. 788 — PHI 318205 — AGRW ID# 23353: %
“Boundary marker of the grave of the society (\textit{thiasos}) of grain-measurers (?) (\textit{sitometrai})” / ὅρος θηκαίων | θιάσου σιτο|με[τρῶν ? — —]. %
See \url{philipharland.com/greco-roman-associations/grave-boundary-marker-of-a-society-of-grain-measurers-ii-ce/}.} \\[0.4mm] %
\noindent{}\textbf{συγκακοπαθέω} / \textit{sugkakopatheō} / “to partake in suffering” (1Tim 1:8, 2:3; \\\makebox[3em][l]{}Nys, \textit{orat}. 23 PG35:1225) \\[0.4mm] %
\noindent{}\textbf{συγκακουχέομαι} / \textit{sugkakoucheomai} / “to endure adversity with” (Heb 11:25; \\\makebox[3em][l]{}Nys, \textit{Cant}. PG6:354) \\[0.4mm] %
\noindent{}\textls[-15]{\textbf{συμμιμητής} / \textit{summimētēs} / “co-imitator” (Phlp 3:17; BasilC, \textit{Reg.Mor.} PG31:848)}\\[0.4mm] %
\noindent{}\textbf{συμπρεσβύτερος} / \textit{sumpresbuteros} / “co-elder” (1Pet 5:1; \textit{Adv.Cataphyrg.} fr. 1) \\[0.4mm]
\noindent{}\textbf{συνομορέω} / \textit{sunomoreō} / “to border on” (Acts 18:7; JChrys, \textit{ActH} PG60:276) \\[0.4mm] %
\noindent{}\textbf{σφυδρόν} / \textit{sphudron} / “ankle”|“foot-arch” (Acts 3:7; Hesychius, \textit{Lex}. Σ 2931)\endnote{%
A late derivation from the Homeric and subsequently commonplace σφυρόν / \textit{sphuron} / “ankle”.} \\[0.4mm] %
\noindent{}\textbf{ταβέρνα} / \textit{taberna} / “tavern” (Acts 28:15; JChrys, \textit{ActH} PG60:374, 376) \\[0.4mm] %
\noindent{}\textls[-15]{\textbf{τυφωνικός} / \textit{tuphōnikos} / “tempestuous” (Acts 27:14; JChrys, \textit{ActH} PG60:368)}\endnote{%
The adjective corresponding to a named entity, Typhonian, has the same spelling, but Acts is apparently the first extant example of the adjective lemma based on a common noun.} \\[0.4mm] %
\noindent{}\textls[-10]{\textbf{ὑπερλίαν} / \textit{huperlian} / “beyond measure” (2Cor 11:5, 12:11; Athan, \textit{Or3Ar.} 24.6) }\\[0.4mm] %
\noindent{}\textbf{ὑπερπερισσῶς} / \textit{huperperissōs} / “beyond all measure” (Mk 7:37; Theophylact, \\\makebox[3em][l]{}\textit{Enarr.} 2.565)\endnote{%
Excepting NT catena, this puts the earliest reception of this lemma in the TLG in the 11th--12th century CE!} \\[0.4mm] %
\noindent{}\textls[-15]{\textbf{φιλοπρωτεύω} / \textit{philoprōteuō} / “to wish first-place” (3Jn 1:9; Pallad, \textit{DialChrys} 131)} \\[0.4mm] %
\noindent{}\textbf{χαλκηδών} / \textit{chalkēdōn} / “chalcedony” (Rev 21:19; Andreas, \\\makebox[3em][l]{}\textit{Apoc.com.} 23.67.21,19d)\endnote{%
The lemma occurs as a named entity noun, Chalcedon, and as an adjective corresponding to that noun, %
in \textit{SibOrac} 3, Herodian and Ps-Herodian, et al, as well as Greek papyri, most notable %
and earliest among these, p.oxy. 11 1380 — TM 63688 — 98--136 CE. %
Besides Revelation, other candidates for the earliest use of the adjective not evoking the named entity are: %
Perinthos-Herakleia 124 — PH167335 — Thrace — Perinthos-Herakleia — Marmara Ereğlisi — 2nd c. CE, %
like Revelation, also in reference to stones. %
It does not seem to re-emerge until the first major commentary on Revelation, %
that by Andreas of Cretensis, in the 6th--7th century CE.} \\[0.4mm] %
\noindent{}\textbf{χαλκολίβανον} / \textit{chalkolibanon} / “fine brass” (Rev 1:15, 2:18; Epiph, \textit{Pan} 51.33)\endnote{%
Holl \& Dummer (GCS 31) 1980:308.} \\[0.4mm] %
\noindent{}\textbf{χρυσοδακτύλιος} / \textit{chrusodaktulios} / “gold-fingered” (Jas 2:2; CyrAl, \\\makebox[3em][l]{}\textit{Jo.com.} Pusey 1872:504)\endnote{Canadian costumery and Scottish voiceover courtesy of Mike Myers.} \\[0.4mm] %
\noindent{}\textbf{χρυσόπρασος} / \textit{chrusoprasos} / “chrysoprase” | “golden-green” (Rev 21:20; \\\makebox[3em][l]{}Oecumenius, \textit{Rev.com.}, Hoskier 1928:240)\endnote{%
There are two dubiously dated hits in the TLG in the so-called \textit{Seniores Alexandrini}, %
citing Ps-Xenophon, \textit{Lithognomon}; see Pitra 1884:342.} \\[0.4mm] % %

\end{small} 

\noindent{}\textit{Summary}: Acts, the Pastorals, 2Cor 8--11, the Catholic Epistles, and Revelation are overrepresented in this list. %
The shared newness/strangeness of their terms and texts reinforce the late 2nd, early 3rd century CE compositional and/or redactional dates for these works %
demonstrated in the main chapters of this book, as well as the priority of the stylistically simpler and more traditional Marcionite NT to its more innovative, proto-Orthodox canonical counterparts. \\[1mm] %

\section{\textit{I Enoch Graecus}}

\noindent{}\textbf{θορυβάζω}|\textbf{θορυβάζομαι} / \textit{thorubadzō}|\textit{thorubadzomai} %
/ “to trouble”|“to be troubled” appears only 13 times in the TLG. %
The first clear, non-Christian use is by the 4th century CE %
grammarian Dositheus (\textit{Ars.gram.} 71). This lemma, incidentally, should not be confused %
with the highly similar θορυβέω / \textit{thorubeō} / “to disturb”|“to raise a clamor”, %
which was common in Attic Greek and subsequently (3501 TLG hits). %

\textit{First Enoch}, specifically the portion known as the \textit{Book of the Watchers} (IEnoch 1--36), %
largely relies on a Ge'ez/Ethiopic manuscript for its full text, %
but portions of it are extant in Greek as well (Black 1970). %
Among the Greek excerpts is the following, in which the rarer lemma appears: % 

\vspace*{3mm} 
 \begin{quote}ἰδοὺ νεφέλαι ἐν τῇ ὁράσει
 ἐκάλουν καὶ ὁμίχλαι με ἐφώνουν, καὶ διαδρομαὶ τῶν ἀστέρων καὶ διαστραπαί
 με κατεσπούδαζον καὶ \textbf{ἐθορύβαζόν} με, καὶ ἄνεμοι ἐν τῇ ὁράσει μου ἐξεπέτασάν με %
 ⁹ καὶ ἐπῆράν με ἄνω καὶ εἰσήνεγκάν με εἰς τὸν οὐρανόν, καὶ εἰσῆλθον
 μέχρις ἤγγισα τείχους οἰκοδομῆς ἐν λίθοις χαλάζης καὶ γλώσσης πυρὸς
 κύκλῳ αὐτῶν· καὶ ἤρξαντο ἐκφοβεῖν με. (\textit{I Enoch} 14.8--9) \\[1mm] %
 Behold, clouds in the vision were calling and mists beckoned me, %
 and the courses of the stars and long-lightning %
 deeply concerned me and \textbf{troubled} me, and winds in the vision of mine unfurled me %
 ⁹ and lifted me up and brought me into the heaven. %
 And I entered until I was near a wall built with stones of hail, and a tongue of fire %
 [was] around them. And they began to terrify me. %
 \end{quote} % 
 \vspace*{3mm}

This portion is extant in the Ethiopic witness, %
and the phrase with the relevant verb is translated, %
“and causing me to desire” (trans. Isaac).\endnote{%
Trans. E. Isaac in OTP1 20.} %
In the footnote, Isaac notes though that “[s]ome think this to be inaccurate”. %
While I do not pretend to competence in Ethiopic, %
the expression as translated may reflect a causative (\textit{hiphil} or \textit{aphel}) verb %
in a Hebrew and/or Aramaic precursor, %
even if the Ethiopic derives from a Greek translation. %
The original may reflect a highly significant Judean word \texthebrew{חמד} / “to covet|desire”, %
typically translated by ἐπιθυμέω in Greek. %
Such an expression puts Greek \textit{1En} 14.8 close to canonical Rom 7:7--9 and Rabbinic thought more generally. %
This quick, initial assessment suggests that, even if %
1 Enoch was translated into Greek prior to the 1st century CE, %
the rare Greek lemma θορυβάζω was unlikely part of that text, %
for which the earliest Greek fragments date to the 4th century CE. %

Whether the rare lemma first appeared in Greek \textit{1En} or canonical Luke is %
debatable, but its Lukan usage was almost %
certainly absent from Marcionʼs \textit{Evangelion}. %
It appears in a rebuke of Martha given by Jesus %
as a guest in the home of Mary and Martha (Lk 10:41): \clearpage

 \begin{quote}Μάρθα Μάρθα, μεριμνᾷς καὶ \textbf{θορυβάζῃ} περὶ πολλά. \\[1mm] %
 Martha, Martha, you are worried and \textbf{troubled} about many things.\end{quote} % 
 \vspace*{2mm} 

\noindent{}If Greek \textit{1En} is the point of origin and influence, %
then Jesus as guest may speak as an Enoch-like apocalyptic traveler and teacher of wisdom. %
Given the association of hospitality and divine/angelic visits %
in traditional Hebrew and Greek thought, %
this line of interpretation does not seem like a stretch. %
If canonical Luke is the point of origin and influence, %
then Enoch imitates Jesus precisely when he is shown hospitality in paradise. %
Either possibility is interesting, %
but either way the evidence pushes toward a late provenance for this neologism. \\[1mm] %

\section{\textit{4 Baruch}}

\noindent{}4 Baruch, also known as the \textit{Paraleipomena of Jeremiah}, was most likely originally written in Greek. %
It has vexed scholars in terms of its original date of composition, with theories varying from the late 1st century CE to the 130s CE.\endnote{%
Jones 2011:148--156, with literature.} %
If the sixty-six year sleep of Abimelech refers back to the destruction of the Jerusalem temple, %
as some scholars surmise, then that brings us to the end of the bar Kochba revolts and %
close to the end of Hadrianʼs reign. %
The irenic focus on remedying relationships between Judeans and Samaritans, as well as Judeans and Romans, favors a context shortly after Hadrian. \\[0.5mm] %

\noindent{}\textbf{ἀλίσγημα} / \textit{alisgēma} / “pollution” appears once in the NT, in Acts 15:20, %
where the Jerusalem council sends out a letter with guidance “to abstain from the \textbf{pollutions} of idols” / %
τοῦ ἀπέχεσθαι τῶν \textbf{ἀλισγημάτων} τῶν εἰδώλων. %
Its earliest appearance and point of origin may well be in \textit{4 Baruch} 7:37, %
which may have inspired Acts at this point: %

\vspace*{1.5mm}
 \begin{quote}Ἱερεμίας δὲ ἄρας τὰ σῦκα διέδωκε τοῖς νοσοῦσι τοῦ λαοῦ, %
 καὶ ἔμεινε διδάσκων αὐτοὺς τοῦ ἀπέχεσθαι ἐκ τῶν \textbf{ἀλισγημάτων} τῶν ἐθνῶν τῆς Βαβυλῶνος. \\[1mm] %
But Jeremiah, taking the figs, distributed to the sick among the people; %
 and he remained, teaching them to abstain from the \textbf{pollutions} of the gentiles of Babylon. \end{quote} % 
 \vspace*{1.5mm}

 \noindent{}\textbf{μισθαποδοσία} / \textit{misthapodosia} / “wage-paying” appears three times in the NT, %
 all in Hebrews (2:2, 10:35, 11:26), %
 two of which are cited explicitly by ClemAlex (\textit{Strom}. 4.16.101.3, 4.16.103.2). %
 Ps-Clementine literature has two examples (\textit{Recog}. 2.21.2--3; \textit{Epit}. 62), %
 but they are not clearly dependent on any of the three references in Hebrews. %
 \textit{4 Baruch} may be its earliest attestation, %
perhaps even its point of origin. %

\vspace*{3mm}
\noindent{}\textbf{μισθαποδότης} / \textit{misthapodotēs} / “wage-payment” appears once in the NT, %
specifically in Hebrews (11:6). Its next secure attestation is by ClemAlex in the %
\textls[-10]{a later book (ca. 203--214 CE) of the \textit{Stromata} (6.9.75.2). %
\textit{4 Baruch} may reflect the earliest attestation, and literary point of origin, for this word as well.} The lemma does not seem to appear %
in Greek papyri, and its only apparent attestation %
in digitized Greek inscriptions is a clearly Christian dedication %
of the 4th century.\endnote{%
SGLIBulg 207 Thracia — Philippopolis (Plovdiv) — 4th c. AD — SEG 14.486 — PH167766. See \url{epigraphy.packhum.org/text/167766}} \\[1mm] %%

\section{\textit{4 Maccabees}}

\noindent{}\textbf{βασανισμός} / \textit{basanismos} / “torture” in the TLG has %
three hits for the 4th--3rd c. BCE playwright, Alexis Comicus. %
However, these three hits reflect different scholarly editions quoting the same, %
solitary, second-hand attestation by the early 3rd c. CE sophist and humorist Athenaeus.\endnote{%
\textit{Deipn}. 1.56 (Kaibel 1887 1:71), quoted as Frag. 290 in Kock 1884:401; Play FIF fr. 23 in Meineke 1840:515; Austin \& Kassel 1991, Play 139, fr. 292.}
According to Athenaeus, Alexis told this joke: %

\vspace*{3mm} 
 \begin{quote}οἶνος ξενικὸς παρῆν· ὁ γὰρ Κορίνθιος \textbf{βασανισμός} ἐστι. \\[1mm] %
 Foreign wine is available. For the Corinthian [wine] is \textbf{torture}.\end{quote} % 
 \vspace*{3mm}

There is ample reason to doubt the authenticity of this %
humorous bit placed among the banter of the \textit{Sophistsʼ Banquet}. %
No other hits for this lemma occur before the 2nd century CE, %
yet a sudden rash begins to appear in \textit{4 Maccabees} and Revelation (see below). %
One wonders whether the joke could have alluded to the famously contentious %
dealings between the sophist Favorinus and the city of Corinth. %

This interpretation must remain tentative, however. %
Alexis was known for his creative word-play. %
Even though Alexis was not apparently a resident of Corinth, %
ruling out the possibility of self-reference, %
the characters in his comedies could have been set in Corinth. %
Euripidesʼ \textit{Medea} and \textit{Alcmaeon} made a backdrop of this famous city. %
Moreover, several cognate lemmas are securely attested quite early: %
\begin{spacing}{0.8}
\begin{small}
\begin{itemize}
\item βασανίζω / \textit{basanidzō} / “to torture” (Heraclitus, 6th--5th BCE)%
\item βασανιστής / \textit{basanistēs} / “torturer” (Antiphon, 5th BCE)%
\item ἀβασάνιστος / \textit{abasanistos} / “not tortured” (Antiphon, 5th BCE)%
\item ἀβασανίστως / \textit{abasanistōs} / “without pain” (Thucydides, 5th BCE)%
\end{itemize}
\end{small}
\end{spacing} 

\noindent{}Other cognate lemmas, however, are first attested quite late: %

\begin{spacing}{0.8}
\begin{small}
\begin{itemize}
\item βεβασανισσμένως / \textit{bebasanissmenōs} / “with severe scrutiny” (Pollux, \\ % 
 \makebox[2em][l]{} ca 178--192 CE; also a favorite of Origenʼs)%
\item βασάνισμα / \textit{basanisma} / “torture” (Theodore Studite, 8th--9th c. CE)%
\end{itemize}
\end{small}
\end{spacing} 

Our focal lemma βασανισμός / \textit{basanismos} / “torture” occurs twice %
in 4 Maccabees (9:6, 11:2) and six times in Revelation %
(9:5 \textit{bis}, 14:11, 18:7, 18:10, 18:15). %
With this term 4Mac 9:6 describes the torture endured by “elders of the Hebrews”, %
and 4Mac 11:2 the torture of the “fifth” of the seven sons of an elderly widow, %
all under the Seleucid ruler Antiochus IV Epiphanes. %
Revelation 9:5 repeatedly uses it in regard to a “five month” %
period of torture of Christian saints by the Roman Empire.\endnote{%
This begs the question, could Revelation interpret the seven Maccabean martyr brothers %
as cryptic references to temporal sequences of salvation-historical events?} %

\vspace*{2mm} 
 \begin{quote}καὶ ἐδόθη αὐτοῖς ἵνα μὴ ἀποκτείνωσιν αὐτούς, %
ἀλλ’ ἵνα βασανισθήσονται μῆνας πέντε· καὶ ὁ \textbf{βασανισμὸς} %
 αὐτῶν ὡς \textbf{βασανισμὸς} σκορπίου, ὅταν παίσῃ ἄνθρωπον. \\[1mm] %
 And it was given to them not to kill them, but rather for them to be tortured five months. %
 And the \textbf{torture} of theirs was like a scorpionʼs \textbf{torture}, %
 when it strikes a human. %
 \end{quote} % 
\vspace*{2mm} 

\noindent{}Subsequent uses of the noun (Rev 14:11, 18:7, 18:10, 18:15) reflect karmic reversals %
whereby the Roman Empire falls divinely subjected to torture. \\[1mm] %

\section{\textit{Acts of Carpus, Papylus, and Agathonice}}

\noindent{}\textbf{ἀπεκδύομαι} / \textit{apekdoumai} / “to strip off oneself” (Col 2:15, 3:9; \\ % 
 \makebox[3em][l]{} \textit{Act.Carp.} 9) \\[0.4mm] %
%
\noindent{}\textbf{καυστηριάζω} / \textit{kaustēriadzō} / “to brand” (1Tim 4:2; \textit{Act.Carp.} 12)\endnote{%
The lemma appears without the initial sigma first in Strabo, \textit{Geog}. 5.1.9, %
Aelius Promotus, \textit{Potency} 61, then Iren, \textit{AH} 1.7.6, 1.20.4} \\[0.4mm] %
\noindent{}\textbf{σινιάζω} / \textit{siniadzō} / “to winnow” (Lk 22:31; \textit{Act.Carp.} 17)\endnote{Ign13, \textit{Smyrn}. 7.1.} %
\\[1mm] %

\section{Aristides Apol.}

\noindent{}\textbf{προσκυνητής} / \textit{proskunētēs} / “worshipper” (Jn 4:23; \textit{Frag}. 2.2)\endnote{%
Besides the \textit{Vita Capri} 7, which closely parallels the text of Jn 4:23, %
the next earliest parallel is in ClemAlex \textit{Protrep}. 9.88.1, 10.96.4, 10.97.3. %
The noun lemma does not appear in the major database of Greek papyri, %
and only seems to appear in Greek inscriptions starting around 259--260 CE, %
namely in IGLSyr 7 4028 — PH245340. %
The adjectival cognate lemma προσκυνητός, ‑ή, ‑όν, appears %
in the 2nd--3rd century CE PH267844 — AGRW 13515 and %
the undated (but likely 3rd century or subsequent) St.Pont. III 20 — PH265863. %
Additional examples of the noun and adjectival lemma appear in post-4th century %
Christian artifacts.} \\[1mm] %

\section{Athenagoras}

\noindent{}\textbf{θεοδίδακτος} / \textit{theodidaktos} / “God-taught” (1Thes 4:9, \textit{Legat.} 11.1, 32.4; %
\\ % 
 \makebox[3em][l]{} \textit{ep.Barn.} 21.6; Theophilus, \textit{Autol.} 2.9; Tatian, \textit{Or.Gr.} 29.2--3) \\[1mm] %

\section{Clement of Alexandria}

\noindent{}\textbf{ἀπέκδυσις} / \textit{apekdusis} / “putting off” (Col 2:11; \textit{Strom}. 3.5.43.3) \\[0.4mm] %
\noindent{}\textbf{ἀρχιτελώνης} / \textit{architelōnēs} / “chief-publican” (Lk 19:2; \textit{Strom}. 4.6.35.2)\endnote{%
Using the term in reference to the Lukan Zacchaeus, and equating him with the evangelist Matthew.} \\[0.4mm] %
\noindent{}\textbf{ἀρχιτρίκλινος} / \textit{architriklinos} / “banquet-chief” (Jn 2:8, 9 (bis); \\\makebox[3em][l]{}\textit{Exc.Theod.} 3.65.1)\endnote{%
Outside of Christian sources, the lemma first appears in Heliodorusʼ %
3rd or 4th century \textit{Aethiopica} 7.27.7, also in reference to wine and a feast!} \\[0.4mm] %
\noindent{}\textbf{ἐθελοθρησκεία} / \textit{ethelothrēskeia} / “voluntary religion” (Col 2:23; \\\makebox[3em][l]{}\textit{Strom}. 3.6.51.3) \\[0.4mm] %
\noindent{}\textbf{εὐπάρεδρος} / \textit{euparedros} / “constantly attending” (1Cor 7:35; \\\makebox[3em][l]{}\textit{Strom}. 4.5.21.2, 4.23.149.2) \\[0.4mm] %
\noindent{}\textbf{εὐπερίστατος} / \textit{euperistatos} / “easily besetting” (Heb 12:1; \textit{Strom}. \\\makebox[3em][l]{}4.16.103.1)\endnote{%
 The next hits are not until the late 4th century with Nyssen, Ps-Macarius, Evagrius, et al.} \\[0.4mm] %
\noindent{}\textbf{λογομαχέω} / \textit{logomacheō} / “to war over words” (2Tim 2:14; \textit{Strom}. \\\makebox[3em][l]{}1.10.49.3)\endnote{%
 The verbal lemma was apparently secondary to, and perhaps derived from, %
 the noun lemma that Pomponius Porphyrion, even while writing in Latin, %
 attributed in Greek to Varro, λογομαχία / \textit{logomachia} / “word-war”. %
 See Chrysippus, \textit{Fragmenta moralia}, fr. 449 in Arnim 1903:109.} \\[0.4mm] %
\noindent{}\textbf{ὀφθαλμοδουλεία} / \textit{ophthalmodouleia} / “eye-slavery” (Eph 6:6, Col 3:22; \\\makebox[3em][l]{}\textit{Strom}. 4.8.65.2)\endnote{%
The lemma is most likely missing from Marcionʼs Lao 6:6 and *Col 3:22.} \\[0.4mm] %
 \noindent{}\textbf{πρεσβυτέριον} / \textit{presbyterion} / “elder-council” (Lk 22:66, Ac 22:5, 1Tim 4:14; %
 \textit{Strom}. 6.13.106.2, 6.13.107.3), %
 at least according to a scan of the top 100 date-sorted results, %
 only appears in Christian rather than non-Christian texts during the 2nd through 4th centuries CE. %
 The early Syriac three letter collection of letters attributed to Ignatius of Antioch %
 does not include this term, but instead only a more generic, %
 less technical reference to πρεσβυτέροις / \textit{presbuterois} / “elders”. %
Ign7 is saturated with the new lemma: %
\textit{Eph} 2.2, 4.1, 20.2, \textit{Magn} 2.1, 13.1, \textit{Tral} 2.2, 7.2, 13.2, \textit{Phld} 4.1, 5.1, 7.1, \textit{Smyr} 8.1, 12.2. %
This bevy of references could cloud the fact that, outside of the NT itself, %
ClemAlex in one of his late-life books of the \textit{Stromateis} %
provides the earliest secure authorial and chronological anchor to the technical term, %
thus locating its emergence between 203 and 214 CE.\endnote{%
For an accessible and succinct summary of scholarship about the pre-Severan (190--202 CE, Books 1--5) and post-Severan (203--214 CE, Books 6--8) %
dates of the \textit{Stromateis}, see Itter 2010:33--34.} %
Clement famously delivered a letter from his student Alexander to Antioch around 211 CE, %
and he may have picked up the term there and/or had a hand %
in the editing of the Ignatian seven letter collection. \\[1mm] %
\noindent{}\textbf{προμεριμνάω} / \textit{preomerimnaō} / “to worry beforehand” (Mk 13:11; %
\\\makebox[3em][l]{}\textit{Strom}. 4.9.70.4) \\[0.4mm] %
\noindent{}\textbf{προσωπολήπτης} / \textit{prosōpolēmtpēs} / “person-respecter” (Ac 10:34; %
\\\makebox[3em][l]{}\textit{Strom}. 6.8.63.5) \\[0.4mm] %
\noindent{}\textbf{συζητητής} / \textit{sudzētētēs} / “co-inquirer” (1Cor 1:20; \textit{Strom}. 5.1.8.1)\endnote{%
Also in Ign7 \textit{Eph} 18.1.} \\[0.4mm] %
\noindent{}\textbf{συζωοποιέω} / \textit{sudzōopoieō} / “to co-make-alive” (Eph 2:5, *Col 2:13; %
\\\makebox[3em][l]{}\textit{Protr}. 2.27.2) \\[0.4mm] %
\noindent{}\textbf{συναρμολογέω} / \textit{sumarmologeō} / “to co-fit” (Eph 2:21, 4:16; \\\makebox[3em][l]{}\textit{Strom}. 6.11.95.2)\endnote{%
This lemma is not attested for Lao 2:21 and should be considered a candidate for removal from the TANZ edition.} \\[0.4mm] %
\noindent{}\textbf{συνεκλεκτός} / \textit{suneklektos} / “co-chosen” (1Pet 5:13; \textit{Fragmenta} 9) \\[0.4mm] %
\noindent{}\textbf{συσταυρόομαι} / \textit{sustauroomai} / “to be co-crucified” (Mt 27:44, Mk 15:32, \\\makebox[3em][l]{}Jn 19:32, Gal 2:19, Rom 6:6; \textit{Strom}. 3.11.75.3) \\[0.4mm] %
\noindent{}\textbf{ὑπερέκεινα} / \textit{huperekeina} / “beyond” (2Cor 10:16; \textit{Strom}. 6.18.164.4)\endnote{Also in Orig, \textit{Cels}. 6.26.10, without direct reference back to 2Cor 10:16, %
perhaps suggesting Origen as the progenitor of the lemma.} \\[0.4mm] %
\noindent{}\textbf{ὑπερεκχύννω} / \textit{huperekchunnō} / “to pour out over” (Lk 6:38; \textit{QDS} 33.5)\endnote{%
The older and commonplace lemma spelling is ὑπερεκχέω / \textit{huperekcheō}, %
attested in Straton (3rd BCE) and the LXX. %
Origen also quotes and confirms with ClemAlex yet another unusual form, here at \textit{Exh.mart.} 10. %
Though I had previously thought this lemma was attested for *Ev, %
but Tertullianʼs \textit{fluentem} does not establish the prefixed and oddly spelled form, %
and \textit{Adamantius} book 5 is likely not reliable here, %
so this should be adjusted accordingly. %
Vinzent already adroitly noted this in \textit{CPCNT} 525, %
“Tert. seems to read the \textit{verbum simplex: pressam ac fluentem}”. %
Vinzent did not propose an alternative reconstruction, though. %
The traditional refrain of the land “flowing” / ῥέουσαν (Jer 39:22 LXX) with milk and honey is as good a candidate as any.} %
\noindent{}\textbf{ὑπερεντυγχάνω} / \textit{huperentugchanō} / “to intercede over” (Rom 8:26; \\\makebox[3em][l]{}\textit{Paed}. 6.47.4)\endnote{%
Origen quotes and obsessively uses yet another \textit{huper}-prefixed lemma (9 distinct TLG hits!).} \\[0.4mm] %

\noindent{}\textit{Summary}: ClemAlex has an affinity for certain, gratuitous prefixes (ἀρχ, εὐ, συ, ὑπερ), %
tendencies our team has documented in the redactions of canonical Luke and the canonical Paulines. %
This makes Clement among the strongest candidates of known writers to have contributed to editing the proto-Orthodox, canonical compositions. %
\\[1mm] %

\section{Clement of Rome}

\noindent{}\textbf{δίψυχος} / \textit{dipsuchos} / “double-minded” (Jas 1:8, 4:8; \textit{1Cor} 11:2, 23:2) \\[1mm] %

\section{\textit{Didache}}

\noindent{}\textbf{διώκτης} / \textit{diōktēs} / “pursuer” (1Tim 1:13, \textit{Didache} 5.2, \textit{ep.Barn.} 20.2)\endnote{%
A misspelled form of this word (δηόκτης \ \textit{dēoktēs}) may appear once in Greek papyri: SEG 30:1625 Cyprus Lythrodontas 3rd CE RDAC (1980) 260,1.} \\[0.4mm] %
\noindent{}\textbf{ἐπιούσιος} / \textit{epiousios} / “daily|necessary” (*Ev 11:3, Mt 6:11; \\\makebox[3em][l]{}\textit{Didache} 8.2)\endnote{%
 This is one of the rare instances where a potential neologism appears in Marcionʼs \textit{Evangelion}. %
 Here it happens to appear in the Lordʼs Prayer, in a phrase commonly attributed to the Q Gospel. %
One wonders whether this neologism was coined by Jesus himself, or one of his early followers.} \\ %
\noindent{}\textbf{κατάθεμα} / \textit{katathema} / “accursed” (Rev 22:3; \textit{Didache} 16.5, \textit{Asc.Isa.} 3.18; \\ % 
 \makebox[3em][l]{} Ps-Clem, \textit{Contest}. 4.3)\endnote{%
 This lemma appears in a Greek inscription, but it dates to the 5th or 6th century: %
 IG II\textsuperscript{2} 13517 — Attica Paiania (Liopesi) — SEG 29.250 — E. Sironen, Inscriptions (1997) 258, 226 — SEMA (2006) 2747 — PH345748.} \\[1mm] %

\section{\textit{Gospel of Peter}}

\noindent{}\textbf{ἐπιφώσκω}|\textbf{ἐπιφωσκέω} / \textit{epiphōskō}|\textit{epiphōskeō} / “to be near dawn” %
\\\makebox[3em][l]{}(Mt 28:1, Lk 23:54; \textit{EvPet} 5; Ps-Clem., \textit{hom}. 3.1.1, 3.58.2) %
\endnote{%
Not to be confused with ἐπιφαύσκω / \textit{epiphauskō} / “to shine out”, found in LXX Job 25:5, 31:26, 41:9, Eph 5:14. %
On the \textit{Gospel of Peter} being composed after, and depending on, Marcionʼs \textit{Evangelion}, %
see Vinzent 2014:272--274; Bilby \textit{SGC} at *Ev 23:32--33, Lk2 23:39--43. %
 being earlier than canonical Luke (and likely also canonical Matthew), %
but not (as Crossan would have it), the \textit{Ur}-Gospel, %
see Bilby 2013:29, 39--43, 336--337.} \\[1mm] % %

\section{\textit{Hermas}}

\noindent{}\textbf{ἀνακαίνωσις} / \textit{anakainōsis} / “renewal” (Rom 12:2, Tit 3:5; \textit{Herm}. 16.9)\endnote{%
Thereafter in ClemAlex, \textit{Strom}. 2.9.41.4, 6.15.131.} \\[0.4mm] %
\textbf{βιβλαρίδιον} / \textit{biblaridion} / “small roll” (Rv 10:2, 9, 10; \textit{Herm}. 5:3, 8:3)\endnote{%
Not to be confused with βιβλίδιον / \textit{biblidion} / “little book”, attested in PLille1.7.7 (3rd c. BCE), %
according to the LSJ entry, and also used by Ps-Demosthenes (\textit{Dionysodorum} 1), Antiphanes, and other ancient Greek authors.} %
\noindent{}\textbf{καρδιογνώστης} / \textit{kardiognōstēs} / “heart-knower” (Acts 1:24, 15:8; %
\\\makebox[3em][l]{}\textit{Herm}. 31.4; Ps-Clem., \textit{Hom}. 10.13.3; \textit{ActPaul} 24; \textit{Vit.Carp.} 23) \\[0.4mm] %%
\noindent{}\textbf{κατάλαλος} / \textit{katalalos} / “slanderer” (Rom 1:30; \textit{Herm}. 65, 73, 103; \textit\\ % 
 \makebox[3em][l]{} \textit{ActAndr} 62) \\[0.4mm] %
\noindent{}\textbf{πολύσπλαγχνος} / \textit{polusplagchnos} / “greatly merciful” (Jas 5:11; %
\\\makebox[3em][l]{}\textit{Herm}. 31.5, 57.4, 60.4) \\[0.4mm] %
\noindent{}\textbf{πρωτοκαθεδρία} / \textit{prōtokathedria} / “first-public-seat” (Mt 23:6, 
\\\makebox[3em][l]{}Mk 12:39, Lk 11:43; \textit{Herm}. 43.12) %
\\[1mm] %

\section{Hippolytus of Rome}

\noindent{}\textls[-15]{\textbf{ἐνταφιασμός} / \textit{entaphiasmos} / “entombment” (Mk 14:8, Jn 12:7; \textit{Cant}. 26.3)} \\ %
\noindent{}\textbf{ῥέδη} / \textit{hredē} / “carriage|wagon” (Rev 18:13; \textit{Antichr}. 41)\endnote{%
The lemma is a loanword from the Latin \textit{raeda}. %
In the TLG results, this lemma appears together with a named entity used %
by the astrologer and geographer Ptolemy in his \textit{Geographia} (6.7.41; Grasshoff \& Stückelberger 2:634--635). %
Ptolemy presents Ῥαίδα (Ῥέδα) / \textit{Hraida} (\textit{Hreda}) as the name of a city, together with its coordinates, %
located in the southern Arabian peninsula, in what is today Yemen. %
After Hippolytusʼ close parallel with Revelation, %
the Latin loanword in Greek does not appear again until John Lydus, \textit{De mensibus} 1.32, in the 6th century CE.} \\[0.4mm] %

\section{Irenaeus}

\noindent{}\textbf{αὐτοκατάκριτος} / \textit{autokatakritos} / “self-condemned” (Titus 3:11, %
\\\makebox[3em][l]{}\textit{AH} 1.9.3, 3.5)\endnote{%
 Origen subsequently and repeatedly comments on Titus 3:11.} \\[0.4mm] %
\noindent{}\textbf{ὀκταήμερος} / \textit{oktaēmeros} / “eighth-day” (Phlp 3:5; \textit{AH} 1.11.2) \\[0.4mm] %
\noindent{}\textbf{συγκοινωνός} / \textit{sugkoinōnos} / “co-partner” (Rom 11:17, 1Cor 9:23, %
\\\makebox[3em][l]{}Phlp 1:7, Rev 1:9; Irenaeus, \textit{AH} fr. 10) \\[0.4mm] %
\noindent{}\textbf{σύσσωμος} / \textit{sussōmos} / “co-bodied” (Eph 3:6; \textit{AH} 1.4.1) \\[1mm] %

\noindent{}\textit{Summary}: though he does not attest to nearly as many NT neologisms as ClemAlex or Origen, %
these examples cluster in the canonical Paulines; %
Irenaeus may well have had a hand in their editing. \clearpage %

\section{Justin Martyr}

\noindent{}\textbf{διακαθαρίζω} / \textit{diakatharizō} / “to purge fully” (Mt 3:12; \textit{Tryph}. 49.3)\endnote{
 Next quoted by Irenaeus, \textit{AH} 1.1.6. The generally accepted form is διακᾰθαίρω, %
 a form that goes back to Pherecydes, Aristophanes, and Plato, %
 and continues well into the Common Era, including Luke 3:17. %
 } \\[0.4mm] %%
\noindent{}\textbf{καταθεματίζω} / \textit{katathematidzō} / “to curse down” (Mt 26:74;\\ % 
 \makebox[3em][l]{}  \textit{Tryph}. 47.4 (\textit{bis}); Irenaeus, \textit{AH} 1.7.3, 1.9.3) \\[1mm] %
\noindent{}\textbf{μοσχοποιέω} / \textit{moschopoieō} / “to calf-make” (Acts 7:41, \textit{Tryph}. 19.5, %
\\\makebox[3em][l]{}102.6, 132.1)\endnote{%
μοσχοποιέω / \textit{moschopoieō} does not appear in any digitized Greek papyri or inscriptions, %
strengthening the case for Justin as the lemmaʼs originator.} \\[0.4mm] %
 \noindent{}\textbf{πρόσχυσις} / \textit{proschusis} / “pouring upon” (Heb 11:28; \textit{2Apol} 12.5)\endnote{%
 Among other texts with possible roots in the 2nd century CE, it also appears in Ps-Clem, \textit{Hom}. 2.44.2.} \\[0.4mm] %
 \noindent{}\textbf{ψευδαπόστολος} / \textit{pseudapostolos} / “false-apostle” (*2Cor 11:13; \\\makebox[3em][l]{}\textit{Tryph}. 35.3)\endnote{%
 Also attested in: a fragment attributed by Eusebius to Heggesippus (Diekamp \& Funk 1913:216); Ps-Clem, \textit{Hom}. 16.21.4; ClemAlex, \textit{Strom}. 3.1.3.4; %
 Ign13, \textit{Philad}. 5.2.} \\[0.4mm] %
 \noindent{}\textbf{ψευδοδιδάσκαλος} / \textit{pseudodidaskalos} / “false-teacher” (2Pet 2:1; \\\makebox[3em][l]{}\textit{Tryph}. 82.1)\endnote{%
 Also attested four times by Origen.} \\[0.4mm] %
 \noindent{}\textbf{ψευδόχριστος} / \textit{pseudochristos} / “false-messiah” (Mt 24:24, Mk 13:22; \\\makebox[3em][l]{}\textit{Tryph}. 35.3, 82.2)\endnote{%
 Attested also by Hippolytus, \textit{Cant}. 20.1, \textit{Dan}. 4.18, and Origen, \textit{Cels}. 2.49.} \\[1mm] %

\noindent{}\textit{Summary}: like Irenaeus, Justin only attests a hand-full of NT neologisms, %
but even this small sample shows only one overlap with the Pauline epistles %
that happens to match Marcionʼs text at this point. %
Justin perhaps borrowed this term from Marcionʼs Pauline letters while rejecting them, %
perhaps because the proto-Orthodox were still in the process of deciding whether to appropriate them or not. %
He may also have coined the term ψευδοδιδάσκαλος / \textit{pseudodidaskalos} / “false-teacher” %
to rebut Marcion, %
and either composed, or more likely inspired 2Pet 2:1 at this point. %
Likewise, ψευδόχριστος / \textit{pseudochristos} / “false-messiah” was likely minted by Justin, %
but in this case inserted into the canonical redactions of two Gospels. %
Justinʼs similarly patterned neologisms point to his role as an editor of canonical Matthew. \\[1mm] %

\section{Origen of Alexandria}

\noindent{}\textbf{ἀγενεαλόγητος} / \textit{agenealogētos} / “without genealogy” (Heb 7:3; \\\makebox[3em][l]{}\textit{Jo.com.} 1.4.21) \\[0.4mm] %
\noindent{}\textbf{ἀποκαραδοκία} / \textit{apokaradokia} / “earnest expectation” (Rom 8:19, \\\makebox[3em][l]{}Phlp 1:20; \textit{Cels}. 5.13, 7.65, 8.5) \\[0.4mm] %
\noindent{}\textbf{ἀποσυνάγωγος} / \textit{aposunagōgos} / “excommunicated” (Jn 9:22, 12:42, \\\makebox[3em][l]{}16:2; \textit{Ps.hom.} 25.3) \\[0.4mm] %
\noindent{}\textbf{βαττολογέω} / \textit{battologeō} / “to prattle” (Mt 6:7; \textit{orat}. 2.2, 19.1, 21.1--2)\endnote{%
Note that several cognate words appear in the \textit{Fabulae} of Aesop. Note also that %
Klinghardt is the only editor to restore this word for Marcionʼs \textit{Evangelion}, %
a position that my stylometric research has found to be highly unlikely.} \\[0.4mm] %
\noindent{}\textbf{γυμνιτεύω} / \textit{gumniteuō} / “to go naked” (1Cor 4:11; \textit{Io.com.} 13.2.12) \\[0.4mm] %
\noindent{}\textbf{διερμηνευτής} / \textit{diermēveutēs} / “master-interpreter” (1Cor 14:28; \\\makebox[3em][l]{}\textit{1Cor.cat.} 25) \\[0.4mm] %
\noindent{}\textbf{ἑδραίωμα} / \textit{hedraiōma} / “support” (1Tim 3:15; \textit{Cels}. 5.33) \\[0.4mm] %
 \noindent{}\textbf{καλοδιδάσκαλος} / \textit{kalodidaskalos} / “well-taught” (Titus 2:3; \\\makebox[3em][l]{}\textit{Jo.com.} 32.12.132, \textit{1Cor fr.} 74) \\[0.4mm] %
\noindent{}\textbf{ὀλιγοπιστία} / \textit{oligopistia} / “mini-trust” (Mt 17:20; \textit{Cels}. 1.49; \\\makebox[3em][l]{}\textit{Mt.com.} 10.19, 11.5--6, 12.6; \textit{ActThom} 65) \\[0.4mm] %
 \noindent{}\textbf{πατρολῷος} / \textit{patrolōos} / “parricide” (1Tim 1:9; \textit{Luc.fr.} 87; \textit{Mat.fr.} 44)\endnote{%
 The Attic and far more commonplace Hellenistic lemma is spelled πατραλοίας / \textit{patraloias}. %
 Only apparently with Origen around the mid-3rd century CE does the new spelling in the Pastorals come into wider use.} \\[0.4mm] %
 \noindent{}\textbf{πειθός} / \textit{peithos} / “persuasive” (1Cor 2:4; \textit{Cels}. 1.62, 6.2, etc.)\endnote{%
 This new, abbreviated spelling for πιθανός / \textit{pithanos} seems to appear first in canonical 1Cor 2:4, whereas %
 Marcionite 1Cor 2:4 uses \textbf{πειθώ} / \textit{peithō} / “persuasion” (as art rather than as Goddess), %
 a usage in evidence as early as Aeschylus, \textit{Fragmenta} Tetralogy 36 play B fr. 394. %
 It is possible that the canonical redactor introduced the new form as a spelling mistake, %
 using πειθοῖ[ς] instead of πειθανοῖς as the dative plural adjective. %
 πειθός as a misspelling of πίθος / \textit{pithos} / “jar” is attested in 201 CE at Oxyrhynchus, psi.13.1328.} \\[0.4mm] %
 \noindent{}\textbf{περικρύβω} / \textit{perikrubō} / “to hide entirely” (Lk 1:24; \textit{Luc.hom.} 6.33)\endnote{%
This derives from περικρύπτω / \textit{perikruptō}, but the traditional spelling remained common in Late Antiquity.} \\[0.4mm] %
\noindent{}\textbf{προαιτιάομαι} / \textit{proaitiaomai} / “to pre-accuse” (Rom 3:9; \textit{Rom.com.} 3.9)\endnote{%
προαιτιάομαι / \textit{proaitiaomai} appears only 20 times in the entire TLG, %
all apparently by Christian authors, and most of which are quotations of Rom 3:9.} \\[0.4mm] % 
\noindent{}\textbf{προσωποληπτέω} / \textit{prosōpolēpteō} / “to person-respect” (Jas 2:9; \\\makebox[3em][l]{}\textit{Prov.cat.} PG17:208) \\[0.4mm] %
\noindent{}\textbf{συμμορφίζομαι} / \textit{summorphidzomai} / “to be co-conformed” (Phlp 3:10; \\\makebox[3em][l]{}\textit{Cels}. 2.69) \\[0.4mm] %
\noindent{}\textbf{τροποφορέω} / \textit{tropophoreō} / “to bear a mood” (Acts 13:18; \textit{Cels}. 4.71)\endnote{%
Derived from the previous lemma spelled τροφοφορέω / \textit{trophophoreō}, evident in LXX Deut 1:31 \textit{bis}, 2Mac 7:27. %
Origen interestingly used the τροπ- / \textit{trop}- spelling when quoting LXX Deut 1:31 in his polemic against Celsus, %
which could be worth considering as a tell of his editorial work on Acts, particularly given his particular affinity for this lemma (19 TLG hits!)} \\[0.4mm] %
\noindent{}\textbf{ὑπερεκπερισσοῦ} / \textit{huperekperissou} / “superabundantly” (Eph 3:20, \\\makebox[3em][l]{}1Th 3:10, 5:13; \textit{Orat}. 6.4) \\[0.4mm] %
\noindent{}\textbf{ὑπερπερισσεύω} / \textit{huperperisseuō} / “to super-abound” (Rom 5:20; \\\makebox[3em][l]{}2Cor 7:4; \textit{Rom.com.} 5.21)\endnote{%
ὑπερπερισσεύω / \textit{huperperisseuō} does not occur prior to canonical NT texts, and is first confirmed %
and received by Origen in \textit{Rom.com.} 5.21. %
Our team restored this in our TANZ edition, following the initial assessment of the lemma in \textit{CPCNT} 526. %
Tertullian attests \textit{superabundaret} in (\textit{Marc}. 5.13.10), %
so this lemma on its face makes sense. However, based on broader diachronic patterns, %
it is less likely than another candidate. ὑπερπλεονάζω / \textit{huperpleonadzō}. %
In \textit{CPCNT}, the latter lemma was indexed, but only for 1Tim 1:14. %
However, it is also clearly used in the \textit{Psalms of Solomon} 5.16 (1st BCE -- 1st CE), Heron Mech., \textit{Pneum}. 1.10 (1st--2nd CE), \textit{Hermas} 34.6, and Valens, \textit{Anthol}. 2.22, 3.9, App 11.1 (152--165 CE). %
Neither lemma is evident in searches of major databases of Greek papyri and inscriptions. %
Given the preponderance of evidence, greater weight should be given to the second lemma %
and thus to an alternative reconstruction: ὑπερεπλεόνασεν / \textit{huperepleonasen}. %
This happens to be the exact form found in 1Tim 1:14 and \textit{Hermas} 34.6, %
both of which make sense as downstream receptions of Marcionite Rom 5:21.} \\[0.4mm] %
\noindent{}\textbf{φραγέλλιον} / \textit{phragellion} / “whip” (Jn 2:15; \textit{Jo.com.} x15 (!); \textit{Ps.hom.} x2)\endnote{Among the canonical Gospels, this word appears only in the frontloaded Johannine temple cleansing and is an NT \textit{hapax}. %
As LSJ notes, it is a loanword from the Lat. \textit{flagellum}, %
and later came to mean the “name of a weight”, citing φ. σιδαρᾶ βʹ Petersen-Luschan Reisen in Lykien No.77a (iv A. D.). The word %
may appear, but has to be partly supplied, in SEG 7:372 — PH320613 — Dura-Europos — 2nd CE. Origenʼs obsession with %
this lemma makes him a prime candidate as the redactor who added it.} %
\noindent{}\textls[-15]{\textbf{φραγελλόω} / \textit{phragelloō} / “to flagellate” (Mt 27:26, Mk 15:15; \\\makebox[3em][l]{}\textit{Frag.EvJo} 121)}\endnote{%
\textit{Test.Benj.} 12.1 is another possible candidate for the earliest use of this lemma and/or earliest reception of an NT usage. The preceding cognate term, over which Origen was apparently obsessed, favors this classification. %
The usage in \textit{Test.Benj.} is dramatic. %
Joseph tells Benjamin about his brothers, that %
“[o]ne of them, stripping me of my tunic, gave me a loin-girdle, %
and flagellating me, said to walk” / εἷς ἐξ αὐτῶν ἀποδύσας με τὸν χιτῶνα ἔδωκέ μοι περίζωμα, %
καὶ φραγελλώσας με εἶπε τρέχειν (De Jonge 1978:168). %
Note that the translator of OTP 2:825, H.C. Kee, either uses a quite different Greek text or badly mistranslates the subject, having one of the Ishmaelites beat Joseph, not one of his brothers. The lemma is also used in \textit{Act.Paul.Thec.} fr. 1 and sec. 21, and also may be partly attested in one Greek papyrus, bkt 9 60 -- TM 63371 -- 100--199 CE.} \\[0.4mm] %
\textbf{χρηστολογία} / \textit{chrēstologia} / “fair-speaking” (Rom 16:18; \textit{Lam.fr.} 56; \\\makebox[3em][l]{}\textit{Prov.exp.} PG17:173) \\[1mm] %

\noindent{}\textit{Summary}: Origen was a prolific commentator on Jewish and Christian scriptures, %
but that does not adequately explain the breadth of neologisms he attests. %
More than merely attesting them, he regularly uses many of them even when not referring directly to NT wording. %
Origen shows an easy familiarity with many of these lemmas, %
even to the level of them being habitual. %
Given his education under ClemAlex, and the NT editorial role that Clementʼs many neologisms suggest for him, %
we also suggest that Origen should be considered as a notable and influential editor of the proto-Orthodox Gospels, Acts, %
and the Pauline letters. %
\\[1mm] %

\section{Papias of Hieropolis}

\noindent{}\textbf{εὐαγγελιστής} / \textit{euangelistēs} / “evangelist” (Acts 21:8, Eph 4:11, %
 \\\makebox[3em][l]{}2Tim 4:5; Ign13, \textit{Ant}. 9.4.1, Papias, \textit{fr}. 2.5, 7.1, 9.1) \\[1mm] %

\section{Polycarp, \textit{Philippians}}

\noindent{}\textbf{ἀντίχριστος} / \textit{antichristos} / “antichrist” (1Jo 2:18, 22, 4:3, 2Jo 1:7; \\ % 
 \makebox[3em][l]{} \textit{Phlp} 7.1)\endnote{%
 The term also occurs in chronologically uncertain texts: \textit{Apoc.Elij.} fr. c; \textit{Apoc.Esdr.} 4.31; Ps-Clem, \textit{hom.} 2.17.5; Ign13, \textit{Ant.} 9.5.} \\[0.4mm] %
%
\noindent{}\textbf{ἀρσενοκοίτης} / \textit{arsenokoitēs} / “man-bedder” (1Cor 6:9, 1Tim 1:10; \\ % 
 \makebox[3em][l]{} \textit{Phlp} 5.3; Ps-Clem, \textit{Recog}. 9.23; \textit{ActJo} 36; Ign13, \textit{Tars}. 7.2) \\[0.4mm] %
\noindent{}\textbf{δίλογος} / \textit{dilogos} / “double-tongued” (1Tim 3:8; \textit{Phlp} 5.2) \\[0.4mm] %
\noindent{}\textbf{προσωποληψία} / \textit{prosōpolēmpsia} / “person-respect” (Rom 2:11, Eph 6:9, %
\\\makebox[3em][l]{}Col 3:25, Jas 2:1; \textit{Phlp} 6.1) \\[0.4mm] %
\noindent{}\textbf{ψευδάδελφος} / \textit{pseudadelphos} / “false-brother” (*Gal 2:4, 2Cor 1:26; \\\makebox[3em][l]{}\textit{Phlp} 6.3)\\[1mm] %

\noindent{}\textit{Summary}: based on this small sample of NT neologisms, %
Polycarp is a possible candidate as one of the redactors of the canonical Pauline epistles, including the Pastorals. %
The solitary neologism tied to the Johannine epistles could reinforce the traditional %
narrative of Polycarp as a receiver of a distinctive and earlier Johannine literary tradition. %
Also, Polycarp is the most likely among known authors to have coined %
and/or first popularized the Greek word for “homosexual”. %
Given the apparent novelty of the lemma “false-brother” in Marcionite Galatians, %
Polycarpʼs use of it may reflect the inversion of the insulting label back upon his famous rival. \\[1mm] %

\section{\textit{Protoevangelium of James}}

\noindent{}\textbf{ἐραυνάω} / \textit{eraunaō} / “to search out” (Jn 5:39, 7:52, Rom 8:27, 1Cor 2:10, \\\makebox[3em][l]{}1Pet 1:11, Rev 2:23, \textit{ProtJas} 39)\endnote{%
 The common spelling from Homer through Late Antiquity is ἐρευνάω / \textit{ereunaō}. LXX Deut 13:14[15] in Ralphsʼs edition has a form with this spelling, %
 ἐραυνήσεις / \textit{eraunēseis}, but there are variants here: ἐκζητήσεις / \textit{ekdzētēseis} and ἐτάσεις / \textit{etaseis}. %
 See Wevers 1978:17, 124. Given the peculiarity of the spelling, its prevalence in the canonical redaction, %
 and the theme of searching the scriptures for Jesus (Jn 5:39), %
 this seems to me a highly likely proto-Orthodox revision to LXX Deut that creates a lexical bridge between NT and OT. %
 The same spelling shift is also evident in the lemma ἐξερευνάω (standard) to ἐξεραυνάω (canonical NT). %
 Perhaps a play on words with the verb “to strike with thunderbolts” / %
 κεραυνόω / \textit{keraunoō} was involved.} \\[0.4mm] % %
\noindent{}\textbf{σιρικός} / \textit{sirikos} / “silken”|“Seric”|“Seresian” (Rev 18:12; \textit{ProtJas} 22)\endnote{%
Variants in spelling and gender, as well as overlapping noun and adjective lemmas make this a bit complex. %
Σηρικός as a masculine noun is first attested three different times in the \textit{Periplus Maris Erythraei} (39, 49, 64), which dates %
“without any doubt to the middle of the first century A.D.”, specifically, “between A.D. 40 and 70”, %
based on direct knowledge of the sequence of Nabatean kings (Casson 2012:7). %
Σηρική as a feminine noun is first attested by Teucer of Babylon (\textit{duodecim signis}, Boll 1908:204). %
σηρικός as an adjective probably occurs first in Strabo, \textit{Geog}. 15.1.20 (early 1st century CE), afterwards in Josephus (\textit{bell}. 7.127), %
Plutarch (\textit{conjug.prae.} 145E; \textit{Pyth.orac.} 396B), and other pre-Hadrianic texts. %
But σιρικός as an adjective apparently first occurs in either Revelation or \textit{ProtJas}, whichever comes first.} %
\\[1mm] %

\section{Ps-Clement}

\noindent{}\textbf{διαπαρατριβή} / \textit{diaparatribē} / “constant wrangling” %
(1Tim 6:5; \\\makebox[3em][l]{}\textit{ep.virg.} 1.8.2) \\[0.4mm] %
\noindent{}\textbf{κερματιστής} / \textit{kermatistēs} / “money-changer” %
(Jn 2:14; \textit{Recog}. 9.24) \\[1mm] %
\noindent{}\textbf{οἰκουργός} / \textit{oikourgos} / “working at home” occurs only one time in the NT (Titus 2:5), %
and that term does not occur again in the entire TLG! %
A closely related term --- οἰκουργέω / \textit{oikourgeō} / “to work from home” --- also occurs only once in the entire TLG, %
specifically in Ps-Clem, \textit{1Cor} 1:3, %
apparently revealing a unique connection between the two. %
It is not until the 7th century CE that another lemma with the same stem occurs, %
yet another TLG \textit{hapax}: οἰκουργία / \textit{oikourgia} / “home-work” (Pisides, \textit{Hexaem.} line 684). %
None of these three lemmas seem to appear in any digitized Greek papyri or inscriptions. %
Another closely related term, just without the medial \textit{gamma}, %
οἰκουρός / \textit{oikouros} / “keeping-house” | “house-keeper”, %
was more traditional (Sappho, Aeschylus, et al) and commonplace (255 TLG hits). \\[1mm] %

\section{\textit{Sibylline Oracles} 5}

\noindent{}The \textit{Sibylline Oracles} were an evolving and eclectic collection. %
For book five in particular, I follow the early Hadrianic date (117 CE) %
reached in the careful analysis by Vargas.\endnote{2022:184.} \\[1mm] % %

\noindent{}\textbf{παμπληθεί} / \textit{pamplēthei} / “altogether” %
is an adverb, and must not be confused with the cognate adjective παμπληθής / \textit{pamplēthēs} / “in multitude” | “very numerous”, %
even though the latter can function adverbially, e.g., παμπληθὲς / \textit{pamplēthēs} / “entirely” (LSJ). %
A lemma search of the major Greek papyri database produces a single hit, dated between 1--50 CE, on the main adverb, %
but it is most likely an adjective: “with multitudinous silver” / ἀργύρωι(*) πανπλήθει(*).\endnote{%
p.oxy 9 1184, letter 6a. See \url{papyri.info/dclp/60175}. %
Greek inscriptions with the adverbial lemma do not seem to be in evidence.} %
Thus the first TLG hit, \textit{Sib.Orac.} 5.99, seems to be the earliest extant hit outside the NT. %
It next appears in Justinʼs \textit{Dial.Trypho} 107.2 (ca 155--167 CE) %
and/or \textit{Partitiones} of [Ps-]Herodian, %
Antonine at the earliest.\endnote{%
Boissonade 1819:255.} The use of the adverb in Luke 23:18, in a portion missing from Marcionʼs \textit{Evangelion}, %
associates the term with Barabbas. This may reflect the military import of \textit{Orac.Sib.} 5 %
(there, in reference to the Parthians), or more likely, Cassius Dio (\textit{H.R.} 75.9.1). %
The usage points to canonical Luke being edited after Hadrian, %
and quite possibly in the 2nd quarter of the 3rd century. \\[0.4mm] %

\section{Tatian}

\noindent{}\textbf{αἱματεκχυσία} / \textit{haimatekchusia} / %
“blood-shedding” (Heb 9:22; \\ % 
 \makebox[3em][l]{} \textit{Or.Gr.} 23.2, 23.5) \\[0.4mm] %
\noindent{}\textbf{κατεξουσιάζω} / \textit{kateksousiadzō} / “to exercise authority over” (Mt 20:25, \\ % 
 \makebox[3em][l]{} Mk 10:42; \textit{Or.Gr.} 15.3, 9) \\[0.4mm] %
 \noindent{}\textbf{συλαγωγέω} / \textit{sulagōgeō} / “to carry off as booty” (Col 2:8; %
 \\\makebox[3em][l]{}\textit{Or.Gr.} 22.2, 4) \\[0.4mm] %

\section{\textit{Testament of Abraham}}%

\noindent{}\textbf{ἀνέλεος} / \textit{aneleos} / “merciless” (Jas 2:13; \textit{TestAbrA} 16)%
\endnote{% 
This work is known in two major recensions, Recension A, %
or the longer recension, and Recension B, %
 or the shorter recension. %
 Scholars are mixed about which recension is earlier, %
 or whether both deviate from an older original. 
  See the summary of the history of scholarship in Alison 2003:12--16. %
  It is quite intriguing in this regard that only recension B has several striking parallels %
  with canonical Luke; see Heide 2014:40. %
  Even while favoring the view that Recension A is generally older than Recension B, %
  Alison details an extensive list of words and phrases in Recension A that show %
  that its Greek idiom is “on the whole later” than its counterpart, %
  and possibly “Christian and/or medieval” (Alison 2013:16--19).%
} \\[0.5mm] %

\section{Theophilus, \textit{Autolycum}}

\noindent{}\textbf{ἀντιμισθία} / \textit{antimisthia} / “recompense” (Rom 1:27, 2Cor 6:13; \\ % 
 \makebox[3em][l]{} \textit{Autol}. 2.9)\endnote{%
 Both Pauline verses are missing from Marcion's \textit{Apostolos}. The word %
 also appears in the insecurely attributed and dated Ps-Clem, \textit{2Cor} 1.3, 1.5, 9.7, 11.6, 15.2, and %
 \textit{ActsJohn} 81, before appearing three times in ClemAlexʼs \textit{Stromateis}.} \\[0.4mm] %
 %
 \noindent{}\textbf{ἐπιχορηγία} / \textit{epichorēgia} / “provision” (Eph 4:16, Phlp 1:19; %
 \\ % 
  \makebox[3em][l]{} \textit{Autol}. 1.6, 2.14)\endnote{%
 Both canonical Pauline uses are missing from Marcionʼs \textit{Apostolos}.} \\[0.5mm] %

\section{Valentinus?}
%
\noindent{}While it is quite speculative at this point, %
Valentinus and/or his students may have composed the early version of the %
Deutero-Paulines that ended up in the Marcionite 10 letter collection. %
His greater familiarity with astrological texts likely explains how several of those previously documented neologisms %
made it into the texts that eventually became canonical. %
By my lights, on one occasion, Valentinus may have even coined a new term of his own. \\[1mm] %

\noindent{}\textbf{ἀποκαταλλάσσω} / \textit{apokatallassō} / “to reconcile back” %
(Lao|Eph 2:16, \\\makebox[3em][l]{} *Col 1:20, 22; Irenaeus, \textit{fr.deperd.} 38, \textit{AH} 5 fr. 15) \\[1mm] %

\section{Summary}

\noindent{}These various, apparently Christian neologisms reveal individual tendencies by authors, %
but also show in the aggregate that early Christians focused their creativity %
on developing new terms not for medicine, biology, agriculture, politics, architecture, or even scribal materials,  %
but instead for speaking and speech-acts, %
internecine religious rivalry, %
financial matters, %
ecclesiastical authority, %
monotheistic piety, internal reflection, and divine insight into others. %

By our count, in the main chapters of this book, %
we have surfaced and explored \st{42} 40 total non-Christian neologism candidates. %
In this chapter we identified over 140 total early Christian neologisms. %
This should put to bed the notion that this book presents an argument against the literary and intellectual %
creativity of early Christians and their writings. Quite the opposite. %
These compositions show considerable literary creativity and novelty. %
Christians invented and/or popularized many new words, but not all of the new words. %

Christian writers were products of their time, picking up and imitating words used by their contemporaries. %
But those contemporaries were not merely the likes of Josephus or other late 1st century authors. %
They were also imitating the leading authors of the Second Sophistic. %
This finding effectively synthesizes the work of several scholars on the coordinated creation and/or editing of the texts of the New Testament %
around the mid-second century,\endnote{%
See esp. Klinghardt 2021; Trobisch 1989, 2000, 2023; Vinzent 2023; Livesey 2024.} and on the character of NT texts as exercises in elite literary education, networks, %
and collaborations.\endnote{See esp. Walsh 2021.} %
In other words, the literary tendencies uncovered in this book reveal a relatively small, yet still significant %
and cogent movement of religious Sophists. %
Would we expect anything else of a post-temple Judahite-Greco-Roman religion? %
